\subsection{Dataset and Splits}
The primary benchmark uses \dataset with the selected machine-learning CSV and a fixed 80/10/10 stratified row-level split over the normalized fifteen-class label space. This protocol follows common \dataset intrusion-detection practice and keeps the main result comparable to prior IDS papers. The training split is partitioned into simulated federated clients; validation and test splits remain global and fixed across all methods.

The repository stores split artifacts with seed 20260531. The main non-IID setting uses $K=10$ clients and Dirichlet label skew with $\alpha=0.5$. We also run an IID sanity partition for implementation validation. A $K=20$ client setting, partial participation, and $\epsilon=1$ are optional sensitivity experiments after the main tables are complete. Raw-file per-file 8/1/1 splits and duplicate-controlled diagnostics are treated as robustness checks, not as the main protocol.

\subsection{Backbone and Edge-Model Choice}
The primary and fixed backbone is Qwen3.5-2B, used with the backbone frozen and only adapters trained. This keeps the experimental story focused on a single recent 2B edge-relevant model rather than a broad model-zoo comparison. Qwen3.5-2B supports multimodal inputs, but our protocol uses only text prompts derived from IIoT telemetry; no image input or vision-specific supervision is used.

All main adapter, privacy, and communication comparisons use this same backbone. Cross-backbone checks with Llama 3.2 1B/3B~\citep{meta2024llama32} or Gemma 3 1B~\citep{gemma3technicalreport} are left as optional future or appendix work if reviewers ask for them.

\subsection{Prompt and Features}
The main prompt uses a compact curated feature schema of approximately 20--35 protocol and traffic-behavior fields. Identifier-like and high-leakage columns are excluded, including timestamps, source/destination IPs, payload-like text, URI fields, referer fields, and label columns. Numeric features are rendered consistently, and any quantile binning must use only training data with saved bin boundaries.

\subsection{Compared Methods}
The method comparison is deliberately staged. The P0 set is the minimum publishable package; the P1 set explains why the full method works; P2 is reserved for optional appendix checks.
\begin{itemize}[leftmargin=*]
  \item \textbf{P0 main comparisons}: prompt-only Qwen3.5-2B, centralized LoRA, local-only \core{}, FedAvg-LoRA, Fed-SB-style fixed-core, \core{} without DP, and full \method{} with client-level DP.
  \item \textbf{P1 ablations}: random basis versus public SVD basis, uniform rank versus spectral rank allocation, global DP noise versus layer-adaptive DP noise, no residual versus local residual, and no shrinkage versus shrinkage.
  \item \textbf{Classical IDS context}: Random Forest, XGBoost or LightGBM, and MLP on numeric/tabular features. These baselines calibrate the tabular task; they are not the main private federated assistant comparison.
  \item \textbf{P2 optional checks}: $K=20$, partial participation, $\epsilon=1$, source/protocol skew, and cross-backbone checks.
\end{itemize}
Fed-SB is the strongest collision-risk baseline and must be treated seriously. We first reproduce the authors' public private-FL SNLI experiment to verify environment and privacy behavior. If faithful adaptation to Qwen3.5-2B is brittle, we report a clearly labeled Fed-SB-style fixed-core baseline in our code path under matched budgets.

\subsection{Fair Comparison Regimes}
Rank alone is not a fair matching variable because different adapter methods have different uploaded coordinates, trainable parameters, initialization costs, and privacy effects. Main comparisons are constrained by:
\begin{itemize}[leftmargin=*]
  \item same Qwen3.5-2B backbone, prompt renderer, split, client partition, rounds, and local epochs;
  \item same privacy budget: identical $(\epsilon,\delta)$ and accountant convention for DP methods;
  \item reported communication and capacity budgets: uploaded MB per client per round and trainable adapter parameters.
\end{itemize}
The main epsilon sweep is $\epsilon\in\{8,4,2\}$ with $\delta=10^{-5}$. $\epsilon=1$ is optional and only reported if training remains informative.

\subsection{Metrics and Statistical Reporting}
The primary utility metric is macro-F1 on the fifteen-class test set. We also report balanced accuracy, per-class recall, rare-attack recall, binary AUROC when applicable, and confusion matrices. Efficiency metrics include trainable parameters, uploaded MB per client per round, total communication, peak VRAM, wall-clock time, inference latency, and prompt-token length. Privacy metrics include $(\epsilon,\delta)$, the RDP order grid, clipping rates, and noise multipliers.

All main tables should report at least three seeds with mean and standard deviation. The main comparison against the Fed-SB-style fixed-core baseline should include paired tests or bootstrap confidence intervals. For rare attacks, per-class recall is mandatory because aggregate macro-F1 may hide security-relevant failures.

\begin{table*}[t]
\centering
\caption{Final Qwen3.5-2B experiment matrix. Results are intentionally left blank until server experiments are complete.}
\label{tab:experiment_matrix}
\begin{tabular}{p{0.13\linewidth}p{0.27\linewidth}p{0.32\linewidth}p{0.20\linewidth}}
\toprule
Priority & Question & Methods & Key metrics \\
\midrule
P0 main & Does \method{} improve client-level DP-FL? & FedAvg-LoRA-DP/DP-LoRA-style, Fed-SB-style-DP, \method{}-DP & Macro-F1, rare recall, MB/round, total MB \\
P0 context & What is the non-DP and central/local ceiling? & Prompt-only, central LoRA, local-only \core{}, FedAvg-LoRA, \core{} non-DP & F1, VRAM, wall-clock \\
P0 classical & How strong is tabular IDS? & RF, XGBoost/LightGBM, MLP & F1, rare recall, latency \\
P1 ablation & Which component recovers DP utility? & A0--A5 \method{} variants at $\epsilon=4$ & F1, clipping rate, noise scale \\
P1 diagnostics & Are rankings robust? & Main methods on duplicate-filtered and raw-file checks & Duplicate rate, F1 shift, ranking stability \\
P2 optional & What changes under stress or external validation? & $K=20$, participation $0.5$, $\epsilon=1$, WUSTL-IIOT-2021 & F1, stability, cost \\
\bottomrule
\end{tabular}
\end{table*}

\begin{table}[t]
\centering
\caption{Primary \dataset{} protocol.}
\label{tab:dataset_protocol}
\begin{tabular}{ll}
\toprule
Item & Setting \\
\midrule
Primary file & ML-EdgeIIoT selected CSV \\
Backbone & Qwen3.5-2B, text-only pathway \\
Task & 15-class closed-set IDS \\
Split & Stratified 80/10/10 row-level \\
Seed & 20260531 \\
FL clients & $K=10$ simulated clients \\
Main non-IID & Dirichlet label skew, $\alpha=0.5$ \\
Main DP sweep & $\epsilon\in\{8,4,2\}$, $\delta=10^{-5}$ \\
Optional sensitivity & IID, $K=20$, partial participation, $\epsilon=1$ \\
Privacy unit & Client-level upload release \\
\bottomrule
\end{tabular}
\end{table}

